\section{A covariance principle on the closed measurement body}
\label{sec:covariance83}

Let $\mathcal P_{d,k}$ denote the complete body of ordered $k$-outcome
measurements on $\C^d$: its elements are Hermitian tuples
$E=(E_1,\ldots,E_k)$ with $E_j\succeq0$ and $\sum_jE_j=I_d$.
Here $d\ge1$, $k\ge2$, and $N\ge1$ are integers.  Zero effects,
changing ranks, and noncommuting supports are allowed throughout
this section.  The operational distance $D_N$ uses unhalved trace
norm, retained references, common adaptive control, and at most $N$
calls with bounded public stopping.  The device consumes its input
and returns only its ordered classical label.

The normalized tangent space and its inner product are
\[
 \mathcal T_{d,k}=\{H_j=H_j^*: \textstyle\sum_jH_j=0\},\qquad
 \langle H,K\rangle=\sum_j\tr(H_jK_j).
\]
All operator inequalities below on tuple spaces refer to this real
Hilbert structure, unless stated otherwise.  For any Hermitian tuple
$K$, not necessarily of zero sum, put
\begin{align}
 S_E(K)&=\sum_jE_jK_j,\label{eq:covariances83}\\
 (\mathcal C_EK)_j
  &=\tfrac12\bigl(E_jK_j+K_jE_j-E_jS_E(K)-S_E(K)^*E_j\bigr).
                                      \label{eq:covarianceoperator83}
\end{align}
The matrix $S_E(K)$ need not be Hermitian.  Keeping its adjoint in
\eqref{eq:covarianceoperator83} is essential.  The operator
$\mathcal C_E$ is polynomial in the effects, annihilates constant
tuples, and has range in $\mathcal T_{d,k}$.  In particular its
restriction to $\mathcal T_{d,k}$ is well defined.

For $\tau=N^{-1}$ set
\begin{align}
 g_{N,E}(H)^2
  &=N\langle H,(\mathcal C_E+\tau\operatorname{Id})^{-1}H\rangle,
                                      \label{eq:covarianceform83}\\
 \mathcal Q_N(E,F)&=g_{N,(E+F)/2}(F-E).
                                      \label{eq:covariancemidpoint83}
\end{align}
The inverse acts on $\mathcal T_{d,k}$.  Positivity proved below
makes it nonsingular even on the spectral boundary.  We do not
identify $\mathcal Q_N$ with an operational distance or assume a
triangle inequality for it.

\begin{theorem}[Closed-body adaptive bound]
\label{thm:closedcovariance83}
For every $E,F\in\mathcal P_{d,k}$,
\begin{equation}\label{eq:closedcovariance83}
 D_N(E,F)\le
 \min\{2,\,2\sqrt{k+1}\,\mathcal Q_N(E,F)\}.
\end{equation}
More generally, every continuously differentiable path $E(t)$ in
$\mathcal P_{d,k}$ satisfies
\begin{equation}\label{eq:covariancepath83}
 D_N(E(0),E(1))\le
 \sqrt{k+1}\int_0^1g_{N,E(t)}(E'(t))\,dt.
\end{equation}
These statements are uniform over all ranks and support patterns.
For $k=2$, $E=(A,I_d-A)$ and $H=(B,-B)$, the restriction is exactly
\begin{equation}\label{eq:binarycovariance83}
 (\mathcal C_EH)_1=A(I_d-A)B+BA(I_d-A),\qquad
 \mathcal Q_N(E,F)^2=2Q_N(A,F_1)^2,
\end{equation}
where $Q_N$ is the binary comparison modulus in
\eqref{eq:matrixmodulus76}.
\end{theorem}

Thus the binary variance operator has a normalized-tuple extension,
not an independently regularized list of effects.  The theorem gives
a global upper certificate.  The two-sided balanced-interior
classification and its entropy theorem remain
Theorems~\ref{thm:multioutcomemetric82} and
\ref{thm:multioutcomeentropy82}; a matching general lower bound for
\eqref{eq:covariancemidpoint83} is a separate question.

\subsection{Positive covariance and horizontal decomposition}

\begin{lemma}[Covariance identities]
\label{lem:covarianceidentities83}
On all Hermitian tuples, $\mathcal C_E$ is self-adjoint and positive
semidefinite, and
\begin{align}
 \langle K,\mathcal C_EK\rangle
 &=\sum_j\tr(E_jK_j^2)-\|S_E(K)\|_{\rm HS}^2\notag\\
 &=\sum_j\|E_j^{1/2}(K_j-S_E(K))\|_{\rm HS}^2.
                                      \label{eq:covariancevariance83}
\end{align}
For $E_t=(1-t)E+tF$, $0\le t\le1$,
\begin{equation}\label{eq:covarianceconcavity83}
 \langle K,[\mathcal C_{E_t}-(1-t)\mathcal C_E-t\mathcal C_F]K\rangle
 =t(1-t)\|S_F(K)-S_E(K)\|_{\rm HS}^2.
\end{equation}
The kernel consists precisely of tuples satisfying
$E_j^{1/2}(K_j-S_E(K))=0$ for every $j$.
Moreover $\mathcal C_E=0$ on $\mathcal T_{d,k}$ if and only if the
$E_j$ are mutually orthogonal projections, with zero projections
permitted.
\end{lemma}
\begin{proof}
Write $A_j=E_j^{1/2}$ and let $A$ be the vertical stack of the $A_j$.
Then $A^*A=I_d$, so $P=I-AA^*$ is an orthogonal projection.  On
complex matrix tuples with real Hilbert--Schmidt inner product,
define
\[
 (T_AB)_j=A_j^*B_j+B_j^*A_j.
\]
Its real adjoint is $(T_A^*K)_j=2A_jK_j$.  Direct multiplication gives
\begin{equation}\label{eq:covariancefactor83}
 \mathcal C_E=\tfrac14 T_APT_A^*.
\end{equation}
This proves self-adjointness, positivity, and
\eqref{eq:covariancevariance83}.  It also proves the kernel claim.
Constant tuples are killed and the displayed formula has zero sum,
so no extra projection onto $\mathcal T_{d,k}$ is implicit.

In the first line of \eqref{eq:covariancevariance83}, the first term
is linear in $E$ and $S_E(K)$ is linear in $E$.  Expanding its squared
norm gives \eqref{eq:covarianceconcavity83}.

If the restriction of $\mathcal C_E$ to $\mathcal T_{d,k}$ vanishes,
subtracting the arithmetic mean from any tuple shows that its form
vanishes on all tuples.  Take $K_j=I_d$ at one index and zero
elsewhere.  Its form is $\tr(E_j-E_j^2)=0$.  Since $0\preceq E_j
\preceq I_d$, this forces $E_j^2=E_j$.  Their sum being $I_d$, they
are mutually orthogonal: each positive term in
$\sum_{i\ne j}E_jE_iE_j=0$ must vanish.  Conversely, for orthogonal
projections the summands $E_jK_j$ have mutually orthogonal ranges,
and $\|\sum_jE_jK_j\|_{\rm HS}^2=\sum_j\|E_jK_j\|_{\rm HS}^2$.
Their covariance form is zero.
\end{proof}

\begin{lemma}[Regularized horizontal energy]
\label{lem:horizontalenergy83}
For $H\in\mathcal T_{d,k}$ and $\tau>0$,
\begin{equation}\label{eq:energymin83}
 \langle H,(\mathcal C_E+\tau\operatorname{Id})^{-1}H\rangle
 =\min_{\substack{A^*B=0,\ R\in\mathcal T_{d,k}\\T_AB+R=H}}
       \left(4\|B\|_{\rm HS}^2+\tau^{-1}\|R\|_{\rm HS}^2\right).
\end{equation}
If $K=(\mathcal C_E+\tau\operatorname{Id})^{-1}H$, a minimizer is
\begin{equation}\label{eq:energyoptimizer83}
 B_j=\tfrac12 A_j(K_j-S_E(K)),\qquad R=\tau K.
\end{equation}
\end{lemma}
\begin{proof}
The constraint $A^*B=0$ means $PB=B$.  Thus
$B=PT_A^*K/4$ gives $T_AB=\mathcal C_EK$ and is feasible with
$R=\tau K$.  Its objective is
$\langle K,\mathcal C_EK\rangle+\tau\|K\|^2=\langle H,K\rangle$.
For any feasible perturbation $(\Delta B,\Delta R)$, its linear
objective term is
\[
 8\langle B,\Delta B\rangle+2\tau^{-1}\langle R,\Delta R\rangle
 =2\langle K,T_A\Delta B+\Delta R\rangle=0.
\]
The remaining terms are nonnegative.  This proves minimality,
including at singular $E$; no inverse of an effect was used.
\end{proof}

\subsection{Adaptive differentiation and boundary passage}

\begin{proof}[Proof of Theorem~\ref{thm:closedcovariance83}]
First suppose every effect at the point of differentiation is
positive definite.  The tangent $H_0=T_AB$ in
\eqref{eq:energyoptimizer83} has zero sum.  For small real $s$,
$E_j+s(H_0)_j$ is positive and the tuple remains normalized.
If $R_j(s)$ is its positive square root, then
\[
 U_j=(B_j-R_j'(0))A_j^{-1}
\]
is anti-Hermitian.  Indeed, multiplication of $U_j+U_j^*$ on the
left and right by $A_j$ gives
$A_jB_j+B_j^*A_j-A_jR_j'(0)-R_j'(0)A_j=0$.
The factors $\exp(sU_j)R_j(s)$ therefore realize the derivative
$B_j$ while preserving the effects exactly.

Use the measurement isometry with an observed label and an
inaccessible duplicate label, together with the factor environment.
At the point under consideration it obeys
\[
 W^*\dot W=\sum_jA_j^*B_j=0,\qquad
 \|\dot W\|_{\rm op}\le\|B\|_{\rm HS}.
\]
Purify any fixed common $N$-slot tester.  In the inner product of
terms differentiated at distinct slots, cancel the common later
isometries; the later differentiated slot contains $W^*\dot W$ or
its adjoint and vanishes.  Thus the observed output derivative in
the direction $H_0$ has trace norm at most
$2\sqrt N\|B\|_{\rm HS}$.  This holds for every reference and
memory dimension.  The environments used for comparison are not
returned to the tester.

The output derivative of the same tester is linear in the channel
tangent.  For the residual tuple $R$, one call on a joint state
$\rho$ has reference blocks $\tr_A[(R_j\otimes I)\rho]$.
Duality against Hermitian contractions bounds the trace norm of the
$j$th block by $\|R_j\|_{\rm op}$.  Inserting the residual at each
slot and using trace-norm contraction hence costs at most
\[
 N\sum_j\|R_j\|_{\rm op}\le N\sqrt{k}\|R\|_{\rm HS}.
\]
This is a differential decomposition, not a proposed physical mixture
of measurements.  With $\tau=N^{-1}$, Cauchy--Schwarz and
\eqref{eq:energymin83} now give
\begin{align*}
 2\sqrt N\|B\|_{\rm HS}+N\sqrt{k}\|R\|_{\rm HS}
 &\le\sqrt{k+1}
       \sqrt{4N\|B\|_{\rm HS}^2+N^2\|R\|_{\rm HS}^2}\\
 &=\sqrt{k+1}\,g_{N,E}(H).
\end{align*}
Integrate for the fixed tester, then take the supremum.  Padding a
publicly stopped experiment with discarded dummy calls makes its
original stopped output a common postprocessing of an $N$-slot
experiment, so the same bound applies.

For a boundary path, use $E^\epsilon(t)=(1-\epsilon)E(t)+\epsilon U$,
where $U_j=I_d/k$ and $0<\epsilon<1$.  The covariance is polynomial
and positive semidefinite, so its regularized inverse is continuous.
Moreover $g_{N,E}(H)\le N\|H\|_{\rm HS}$, uniformly on the closed
body.  Dominated convergence applies to the integrals.  The hybrid
bound $D_N(E,F)\le N\sum_j\|E_j-F_j\|_{\rm op}$ gives convergence
of the endpoint distances.  This proves \eqref{eq:covariancepath83}
on the entire body.

For the segment $E(t)=(1-t)E+tF$, put $M=(E+F)/2$ and $H=F-E$.
Concavity and positivity imply, for $0<t\le1/2$,
\[
 \mathcal C_{E(t)}+\tau\operatorname{Id}
 \succeq2t(\mathcal C_M+\tau\operatorname{Id}).
\]
Inverse order bounds $g_{N,E(t)}(H)$ by $(2t)^{-1/2}g_{N,M}(H)$;
use $2(1-t)$ on the second half.  The two scalar integrals sum to
$2$, proving \eqref{eq:closedcovariance83}.  Finally, substituting
$K=(B,-B)$ into \eqref{eq:covarianceoperator83} gives
\eqref{eq:binarycovariance83}.  The factor two in the quadratic form
comes from the two equal Hilbert--Schmidt contributions.
\end{proof}

\subsection{Classical postprocessing}

\begin{proposition}[Data processing of horizontal energy]
\label{prop:covarianceprocessing83}
Let $T=(T_{aj})$ be a column-stochastic classical channel and
$F_a=\sum_jT_{aj}E_j$.  Write $(T^*L)_j=\sum_aT_{aj}L_a$.
Then, on all Hermitian tuples,
\begin{equation}\label{eq:covarianceprocessing83}
 \langle L,\mathcal C_FL\rangle
 \ge\langle T^*L,\mathcal C_E(T^*L)\rangle.
\end{equation}
Consequently the extended quadratic energy
\[
 \mathcal I_E(H)=\sup_K
       \{2\langle H,K\rangle-\langle K,\mathcal C_EK\rangle\}
\]
satisfies $\mathcal I_F(TH)\le\mathcal I_E(H)$.  The value may be
infinite at a boundary; when finite it is
$\langle H,\mathcal C_E^\dagger H\rangle$.
\end{proposition}
\begin{proof}
The two $S$ matrices coincide.  The difference of the covariance
forms is therefore
\[
 \sum_j\tr\!\left[E_j\left(\sum_aT_{aj}L_a^2
                   -(\sum_aT_{aj}L_a)^2\right)\right]\ge0.
\]
The bracket equals the positive sum
$\tfrac12\sum_{a,b}T_{aj}T_{bj}(L_a-L_b)^2$.
Substitution in the variational definition and enlargement from
$K=T^*L$ to all $K$ prove the energy inequality.  The pseudoinverse
statement follows by diagonalizing the real self-adjoint covariance;
a component along its kernel makes the supremum infinite.
\end{proof}

\subsection{A uniform noise crossover, sharp in horizon and noise}

\begin{theorem}[Regularization of arbitrary support patterns]
\label{thm:noisecrossover83}
For $0\le\epsilon\le1$ put $E^\epsilon=(1-\epsilon)E+\epsilon U$.
On $\mathcal T_{d,k}$,
\begin{equation}\label{eq:covariancenoisefloor83}
 \mathcal C_{E^\epsilon}\succeq
       (1-\epsilon)\mathcal C_E+\frac\epsilon k\operatorname{Id}.
\end{equation}
In particular, for arbitrary $E,F\in\mathcal P_{d,k}$,
\begin{equation}\label{eq:noisecrossoverupper83}
 D_N(E^\epsilon,F^\epsilon)\le
 \min\left\{2,
 \frac{2\sqrt{k+1}(1-\epsilon)N}
      {\sqrt{1+N\epsilon/k}}\|F-E\|_{\rm HS}\right\}.
\end{equation}
The dependence on $N$ and $\epsilon$ is attained, up to constants
depending only on $k$, by a two-dimensional family.  Namely, let
$P=|0\rangle\langle0|$, $P_\theta=|u_\theta\rangle\langle u_\theta|$,
$u_\theta=\cos\theta|0\rangle+\sin\theta|1\rangle$, and
\[
 E=(P,I_2-P,0,\ldots,0),\qquad
 F=(P_\theta,I_2-P_\theta,0,\ldots,0).
\]
For $0\le\theta\le\pi/2$, $0\le\epsilon\le1/2$, and
$z=(1-\epsilon)N\sin\theta/\sqrt{1+N\epsilon}$,
\begin{equation}\label{eq:noisecrossoverlower83}
 \frac{\min\{1,z\}}{16384}
 \le D_N^{\rm na}(E^\epsilon,F^\epsilon)
 \le D_N(E^\epsilon,F^\epsilon)
 \le4\sqrt{k(k+1)}\min\{1,z\}.
\end{equation}
\end{theorem}
\begin{proof}
For $K\in\mathcal T_{d,k}$, $S_U(K)=0$ and
$\mathcal C_UK=K/k$.  Concavity proves
\eqref{eq:covariancenoisefloor83}.  Apply it at the midpoint and use
Theorem~\ref{thm:closedcovariance83} to obtain
\eqref{eq:noisecrossoverupper83}.

For the displayed family, keep labels one and two and send every
other label to an independent fair bit.  This common postprocessing
produces the binary effects
$B=\beta P+\epsilon I_2/2$ and
$B_\theta=\beta P_\theta+\epsilon I_2/2$, where $\beta=1-\epsilon$.
The equality holds on subnormalized reference states.  It therefore
preserves the nonadaptive lower witnesses as well as channel
postprocessing.  The midpoint eigenvalues are
$(1\pm\beta\cos\theta)/2$, and direct substitution in
\eqref{eq:matrixcoordinates76} gives
\[
 Q_N(B,B_\theta)^2
 =\frac{4N^2\beta^2\sin^2\theta}
             {2+N(1-\beta^2\cos^2\theta)}
 \ge\frac{4z^2}{2+z^2}.
\]
Here $1-\beta^2\le2\epsilon$ and $N\ge1$ were used.
It follows that $\min\{1,Q_N(B,B_\theta)\}\ge\min\{1,z\}$.
The binary lower theorem in dimension two gives the first inequality
of \eqref{eq:noisecrossoverlower83}.  For the upper,
$\|F-E\|_{\rm HS}=2\sin\theta$ and
$(1+N\epsilon)/(1+N\epsilon/k)\le k$.
Combine \eqref{eq:noisecrossoverupper83} with $D_N\le2$.
\end{proof}

The family passes from projections and zero effects at $\epsilon=0$
to strictly positive effects at every $\epsilon>0$.  Thus a bound
uniform over the closed body must accommodate both the $N$ scale
and the $\sqrt{N/\epsilon}$ scale.  This conclusion does not assert a
matching entropy law on all multi-outcome boundary strata.

\subsection{Exact evaluation without effect square roots}

\begin{proposition}[Polynomial-size rational certificate]
\label{prop:covarianceexact83}
For Gaussian-rational $E,F\in\mathcal P_{d,k}$ and integer $N\ge1$,
$\mathcal Q_N(E,F)^2$ is rational and can be evaluated by exact
linear algebra on a positive definite real system of dimension
$p=(k-1)d^2$.  The bit complexity is polynomial in $d,k$, the total
input bit length, and $\log N$.  Legality is decidable by exact
positive-semidefinite tests.  No effect square root, ODE integration,
rank decision by tolerance, or enumeration of a dictionary is required.
\end{proposition}
\begin{proof}
Use the rational Hermitian matrix basis consisting of diagonal
units and real and imaginary off-diagonal pairs.  For each of the
first $k-1$ slots, insert one matrix basis element there and its
negative in the last slot.  These tuples $B_a$ form a rational basis
of $\mathcal T_{d,k}$; they need not be orthonormal.  Form
\[
 G_{ab}=\langle B_a,B_b\rangle,\quad
 L_{ab}=\langle B_a,\mathcal C_{(E+F)/2}B_b\rangle,\quad
 b_a=\langle B_a,F-E\rangle.
\]
All entries are rational, $G$ is positive definite, and $L\succeq0$.
Solve $(L+N^{-1}G)x=b$.  Then
$\mathcal Q_N^2=N b^{\mathsf T}x$.  Omitting the Gram matrix would
compute a different regularization.  Fraction-free elimination
performs polynomially many rational arithmetic operations; clearing
input denominators and applying determinant bounds gives polynomial
bit length for all minors and for the output.  The same reasoning
applies to exact symmetric elimination for each effect's positivity.
The complexity statement concerns this certificate evaluator, not
the exhaustive public code or collective-readout synthesis.
\end{proof}
